\documentclass[11pt]{article}
\usepackage[a4paper,margin=25mm]{geometry}
\usepackage{amsmath,amssymb,amsthm,mathtools,longtable,booktabs,hyperref}
\usepackage{xurl}
\setlength{\emergencystretch}{2em}
\newtheorem{theorem}{Theorem}
\newtheorem{lemma}[theorem]{Lemma}
\newtheorem{proposition}[theorem]{Proposition}
\newtheorem{remark}[theorem]{Remark}
\newcommand{\Z}{\mathbb Z}
\newcommand{\Q}{\mathbb Q}
\newcommand{\Np}{\mathbb Z_{>0}}
\newcommand{\tri}[1]{\mathcal T(#1)}
\title{A conditional quartic checkpoint from an emended Jones system\\with twelve integrality inputs}
\author{\(\mathbb Z\mathbb Q\) Diophantine Project}
\date{23 November 2026}
\begin{document}
\maketitle
\begin{abstract}
Starting from an explicitly emended Jones-type system, we construct
a rational quadratic system with 255 coordinates and 223 equations, together
with eleven ordinary integer-input requirements and one reciprocal
integer-input requirement. A quadratic definition of the integers in the
rationals with \(m\) total coordinates and \(q_0\) equations therefore gives
a quartic undecidability transfer, conditional on the source-universality
hypothesis $\mathcal U_{\rm em}$, in \(12m+243\) rational unknowns, using
\(223+12q_0\) quadratic squares. The corresponding baseline bounds are
2199, 2391 and 915 for the project's 163-, 179- and 56-coordinate
integrality interfaces. Eliminating one further product register gives the general bound
$12m+242$. Direct optimization of the actual project blocks gives
$12m+239$, namely $2195$, $2387$ and $911$, with $1863$, $2031$
and $771$ quadratic residuals. All hypotheses of the selected integrality
interface remain in force. The coefficient identities and equivalence with
the displayed source are proved. The universality premise is unresolved:
the original 1980 and author-written 1999 displays contain different errors.

The proof uses geometric-quotient integrality, a reciprocal homogeneous
chart, automatic positivity from a cofinite restriction of the instance
parameter, automatic positivity of a Pell congruence quotient, and powers
of a Pell unit to recover positive congruence witnesses. A fixed signed
addition chain computes the large power in Jones's source with 169
quadratic rows. Every step is reversible at the level of existence of
solutions. The accompanying exact checker reconstructs the quadratic core
and its counts.
\end{abstract}

\section{Statement and source scope}
A quadratic integrality interface of size \((m,q_0)\) is a list
\[
 Q_1(v,\mathbf w),\ldots,Q_{q_0}(v,\mathbf w)
 \in\Z[v,w_1,\ldots,w_{m-1}]
\]
of polynomials of total degree at most two such that, for every \(v\in\Q\),
\begin{equation}\label{integer-interface}
 v\in\Z\quad\Longleftrightarrow\quad
 \exists\mathbf w\in\Q^{m-1}\quad
 Q_1(v,\mathbf w)=\cdots=Q_{q_0}(v,\mathbf w)=0.
\end{equation}
Every hypothesis needed to establish \eqref{integer-interface} is part of
the assumptions of a transfer using that interface.

\begin{theorem}\label{main}
Assume \eqref{integer-interface} and Hypothesis $\mathcal U_{\rm em}$
stated below. There is a family of integer-coefficient
polynomials of degree exactly four in
\[
 \boxed{12m+243}
\]
rational unknowns whose solvability is undecidable. Each polynomial is a
sum of \(223+12q_0\) squares of quadratic residuals. The family is indexed
by a positive integer instance parameter, which occurs in its
coefficients and is not counted as an existential unknown.
\end{theorem}

This is a reduction theorem. It does not establish an integrality
interface independently. In particular, ``no additional \(abc\)
hypothesis'' for a project interface does not mean that its other retained
arithmetic premises have been independently discharged.

The system is assembled from Jones \cite{Jones80}, Theorem~3,
p.~861, and Matiyasevich's author-written account \cite{Mat99}, p.~191.
The first prints Greek $\alpha$ in a coefficient where the second prints
Latin $a$; the second has a separate impossible Pell row. Our displayed
system uses the Latin-$a$ form and the correct Pell row. The detailed
source comparison and the exact unverified universality hypothesis are
in Section~\ref{sec:source-reconciliation}.

We prove the complete quadratic compilation, including reconstruction of
positive integer witnesses. We do not infer universality from an
emendation, and none of the numerical quartic bounds here is registered
as a verified undecidability result without $\mathcal U_{\rm em}$.
The alleged 57-variable/47-quadratic theorem has not been independently
reproved. A reversible 62-variable/50-row to 57-variable/45-row reduction
of the literal Mishchenko--Treier transcription is recorded separately;
its universality also remains unresolved.

\section{The precise emended source system}
Fix positive integer parameters \(z,u,y\). Under Hypothesis
$\mathcal U_{\rm em}$ they can be chosen so that the solvable positive
inputs form a computably enumerable nonrecursive set \(W\). Put
\[
 L=5^{60}=867361737988403547205962240695953369140625.
\]
The emended system has the following 28 positive integer unknowns:
\[
 a,b,c,d,e,F,g,h,i,j,k,\ell,m,n,o,p,q,r,s,t,w,
 \alpha,\gamma,\eta,\theta,\lambda,\tau,\varphi.
\]
We write the source root \(f\) as \(F\) to distinguish it from the final
polynomial. The eighteen source equations are
\begin{align}
 e\ell g^2+\alpha&=(b-xy)q^2,\label{J1}\\
 q&=b^L,\label{J2}\\
 \lambda+q^4&=1+\lambda b^5,\label{J3}\\
 \theta+2z&=b^5,\label{J4}\\
 \ell&=u+t\theta,\qquad e=y+m\theta,\label{J5}\\
 n&=q^{16},\label{J7}\\
 r&=\bigl[g+eq^3+\ell q^5+
     \{2(e-z\lambda)(1+xb^5+g)^4+\lambda b^5+
          \lambda b^5q^4\}q^4\bigr](n^2-n)\notag\\
 &\quad+\bigl[q^3-b\ell+\ell+\theta\lambda q^3+
                    (b^5-2)q^5\bigr](n^2-1),\label{J8}\\
 p&=2ws^2r^2n^2,\label{J9}\\
 p^2k^2-k^2+1&=\tau^2,\label{J10}\\
 4(c-ksn^2)^2+\eta&=k^2,\label{J11}\\
 k&=r+1+h(p-1),\label{J12}\\
 a&=(wn^2+1)rsn^2,\label{J13}\\
 c&=2r+1+\varphi,\label{J14}\\
 d&=bw+ca-2c+4a\gamma-5\gamma,\label{J15}\\
 d^2&=(a^2-1)c^2+1,\label{J16}\\
 F^2&=(a^2-1)i^2c^4+1,\label{J17}\\
 (d+oF)^2&=\bigl((a+F^2(d^2-a))^2-1\bigr)
                      (2r+1+jc)^2+1.\label{J18}
\end{align}
The line \eqref{J5} contains two equations. In particular, the exponent
of \(n\) in \eqref{J9} is \emph{two}; the last norm coefficient is a
square minus one; and \eqref{J17} contains \(i^2c^4\), with no additive
one inside that product. These are decisive source features.

We will restrict the instance to
\begin{equation}\label{tail}
 x>2z.
\end{equation}
This restriction discards finitely many instances of the one fixed
nonrecursive set. It preserves undecidability: a decision procedure for
the remaining tail, combined with a finite hardcoded table for the
omitted instances, would decide the entire set. No variable or equation
is required to impose \eqref{tail}, since \(x\) and \(z\) are
coefficients of each instance polynomial.

\section{Three elementary lemmas}
\begin{lemma}[Rational roots of monic equations]\label{monic}
If a rational number satisfies a monic polynomial with integer
coefficients, it is an integer. In particular, a rational number whose
square is an integer is an integer.
\end{lemma}
\begin{proof}
Write the number as \(A/B\) in lowest terms with \(B>0\), substitute
in a monic equation of degree \(d\), and multiply by \(B^d\).
Reduction modulo \(B\) gives \(B\mid A^d\), so \(B=1\).
\end{proof}

\begin{lemma}[Three triangular witnesses]\label{triangular}
For an integer \(v\),
\[
 v>0\quad\Longleftrightarrow\quad
 \exists t_1,t_2,t_3\in\Q\quad
 2v-2=\sum_{j=1}^3t_j(t_j+1).
\]
Moreover, the displayed equality implies \(v>0\) even if \(v\) is
initially only rational.
\end{lemma}
\begin{proof}
Each \(t(t+1)/2\) is at least \(-1/8\), so the equality implies
\(v\ge5/8>0\). Conversely, Gauss's theorem that every nonnegative
integer is a sum of three triangular numbers applies to \(v-1\).
\end{proof}

\begin{lemma}[A reciprocal integer input without an extra coordinate]
\label{reciprocal}
Suppose an integer input \(i\) of \eqref{integer-interface} occurs in a
quadratic construction through \(E=iC\), with \(C\ne0\). Replace the
coordinate \(i\) by \(v\), replace this equality by \(vE=C\), and
replace each integrality residual by
\[
 \widehat Q_j(v,\mathbf w')
    =v^2Q_j(1/v,\mathbf w'/v).
\]
These are polynomials of degree at most two. Their rational solutions
are equivalent to the original solutions with \(i\ne0\), by
\(i=1/v\) and \(\mathbf w=\mathbf w'/v\). The variable count is
unchanged.
\end{lemma}
\begin{proof}
Every monomial of \(Q_j\) has degree at most two, so multiplication
by \(v^2\) clears its denominator. The equation \(vE=C\ne0\)
forces \(v\ne0\). Division by \(v^2\) then recovers exactly
\eqref{integer-interface} at \(i=1/v\). Conversely, multiply each
old auxiliary witness by \(v=1/i\). No point at \(v=0\) survives.
\end{proof}

\section{A quadratic circuit for the source}
The twelve initially free coordinates are
\begin{equation}\label{free12}
 b,g,t,m,w,s,h,\varphi,\gamma,o,j,v.
\end{equation}
We attach ordinary integrality requirements to the first eleven, and a
reciprocal integrality requirement \(1/v\in\Z\) to the last. Three
further free coordinates are \(\lambda,F,\tau\). All other core
coordinates will be outputs of the following quadratic relations, except
for the nonzero reciprocal chart already described.

\subsection{The fixed power}
Put \(P_1=b\). For an addition step \(e=a+b\), introduce \(P_e\)
and the equation
\[
 P_e-P_aP_b=0.
\]
For a subtraction step \(e=a-b>0\), introduce \(P_e\) and the equation
\[
 P_eP_b-P_a=0.
\]
Provided the original \(b\ne0\), induction shows that every coordinate
is exactly \(P_e=b^e\), including all quotient steps. This is an exact
power identity, not a general inference that dividing integers preserves
integrality. The initial addition steps establish $P_5=b^5$ before any quotient
is taken. Therefore the guard $P_5-2z>0$ below forces $b\ne0$
without a circular appeal to a quotient step.

The following signed binary expansion determines our chain:
\[
 L=\sum_{(a,d)\in\mathcal D} d\,2^a,
\]
where
\begin{align*}
\mathcal D=\{&(0,-15),(5,13),(10,3),(16,-7),(21,-13),\\
 &(27,9),(32,9),(38,11),(45,1),(50,9),\\
 &(57,-5),(63,-5),(69,1),(76,5),(82,9),\\
 &(89,-15),(94,7),(99,15),(104,1),(109,-13),\\
 &(114,-3),(119,5),(124,15),(130,-3),(137,5)\}.
\end{align*}
First compute \(P_2\), followed by
\(P_3,P_5,P_7,P_9,P_{11},P_{13},P_{15}\), using addition of the
exponent 2. This costs eight coordinates. Start the running value at
\(P_5\), the leading digit. Reading the remaining positions from 136
down to 0, double its exponent at every position, and add or subtract
the indicated odd exponent when the digit is nonzero. There are 137
doublings and 24 digit updates, nine of which are subtractions.
Every intermediate exponent is positive. Thus the total is
\[
 8+137+24=169
\]
new coordinates and quadratic rows. Set \(B=P_5\) and \(q=P_L\).
Both are already among these coordinates. The companion checker verifies
the signed expansion and every step with exact integers.

\subsection{The remaining rows}
Introduce the power coordinates
\[
 q_2=q^2,\quad q_3=qq_2,\quad q_4=q_2^2,\quad q_5=qq_4,
 \quad q_8=q_4^2,\quad n=q_8^2,\quad n_2=n^2.
\]
The names are coordinates, not unexpanded high-degree expressions.
Define the affine expressions
\[
 \theta=B-2z,\qquad c_0=1+xB+g.
\]
Retain the quadratic relation
\[
 \lambda(B-1)=q_4-1.
\]
Introduce
\begin{align*}
 \ell&=u+t\theta,&e&=y+m\theta,\\
 c_{02}&=c_0^2,&c_{04}&=c_{02}^2,&V&=c_{04}q_4,\\
 S&=g+eq_3+\ell q_5+2(e-z\lambda)V
                      +(\lambda+q_4-1)(q_4+q_8),\\
 T&=(1-b)\ell+[q_4+(1-2z)\lambda]q_3+(B-2)q_5,\\
 r&=S(n_2-n)+T(n_2-1),\\
 e\ell_*&=e\ell,&g_2&=g^2.
\end{align*}
Here \(e\ell_*\) is one coordinate naming the product \(e\ell\);
in the code it is named \texttt{el}. Define the quadratic expression
\[
 \alpha=(b-xy)q_2-e\ell_*g_2.
\]
The regroupings of \(S,T\) use only
\(\lambda B=\lambda+q_4-1\), and hence recover exactly the two
brackets in \eqref{J8}.

Continue with
\begin{align*}
 s_r&=sr,&s_{r2}&=s_r^2,&U&=wn_2,&p&=2Us_{r2},\\
 Y&=sn_2,&M&=rY,&a&=(U+1)M,\\
 k&=r+1+h(p-1),&p_k&=pk,\\
 \tau^2-p_k^2+k^2-1&=0.
\end{align*}
Let \(c=2r+1+\varphi\) be an affine expression rather than a separate
coordinate. Introduce \(Y_K=Yk\), and define
\[
 \eta=k^2-4(c-Y_K)^2.
\]
Finally introduce
\begin{align*}
 d&=bw+ca-2c+4a\gamma-5\gamma,&a_c&=ac,&C_2&=c^2,&d_2&=d^2,\\
 d_2-a_c^2+C_2-1&=0,\\
 vE&=C_2,&A_E&=aE,&F_2&=F^2,\\
 F_2-A_E^2+E^2-1&=0,\\
 G&=a+F_2(d_2-a),&H&=2r+1+jc,&I&=d+oF,&G_H&=GH,\\
 I^2-G_H^2+H^2-1&=0.
\end{align*}
Every displayed residual is quadratic. In particular, the final line is
exactly \(I^2-(G^2-1)H^2=1\).

\subsection{The eleven positivity requirements}
Use Lemma~\ref{triangular}, with a separate triple for each entry, to
impose
\begin{equation}\label{guards}
 g,t,m,w,s,\varphi,\gamma,I,\theta,\alpha,\eta>0.
\end{equation}
The expressions \(\theta,\alpha,\eta\) have degree at most two, so
substitution into their three-triangular residuals preserves quadratic
degree. No separate coordinate for any of them is needed. We will prove
that the omitted source signs, including those of \(r,h,i,j,o,F\), can
be recovered without further positivity coordinates.

\section{Integer reconstruction and automatic positivity}
\begin{lemma}[The geometric quotient]\label{lambda}
Suppose \(b\in\Z\), \(B=b^5>2z>0\), \(q=b^L\), and
\(\lambda(B-1)=q^4-1\). Then \(b\ge2\), and \(\lambda\) is a
positive integer.
\end{lemma}
\begin{proof}
The inequality gives \(b\ge2\) and \(B>1\). Put
\(J=4\cdot5^{59}\). Then \(q^4=B^J\), and the unique rational
value satisfying the equation is
\[
 \lambda=1+B+\cdots+B^{J-1}\in\Np.
\]
\end{proof}

\begin{lemma}[A cofinite restriction removes the sign of \(r\)]
\label{rpositive}
For \(x>2z\), the integer inputs, circuit equations and positivity
conditions in \eqref{guards} already force \(r>0\), without a
positivity wrapper for \(r\).
\end{lemma}
\begin{proof}
The power equations and \(\theta>0\) give \(b\ge2\), and
Lemma~\ref{lambda} gives \(\lambda>0\). Since \(m,t,u,y\ge1\),
\(e,\ell\ge\theta+1>0\). The condition \(\alpha>0\) gives
\[
 e\ell g^2<(b-xy)q^2,
\]
so \(b>xy\ge x>2z\). In particular \(2z\le b-1\), and
\[
 e,\ell\ge b^5-b+2\ge b^4.
\]
Therefore \(g^2<q^2/b^7\), whence \(g<q/b^3\). Also
\(x\le b-1\), and \(L\ge8\), so
\[
 0<c_0=1+xb^5+g<b^6+q/b^3
       \le q(1/b^2+1/b^3)\le 3q/8<q.
\]
The original expression for the first bracket is
\[
 S=g+eq^3+\ell q^5+
 \bigl[2ec_0^4+\lambda\{B(1+q^4)-2zc_0^4\}\bigr]q^4.
\]
Since \(B>2z\) and \(q^4>c_0^4\), all bracketed terms are positive,
so \(S>0\).

Similarly, the original expression for the second bracket is
\[
 T=q^3-(b-1)\ell+\theta\lambda q^3+(B-2)q^5.
\]
The same \(\alpha\) inequality gives \(\ell<bq^2\). Since
\((B-2)q^3>b(b-1)\) for \(b\ge2\) and \(q\ge b\), we obtain
\(T>0\). Finally \(n=q^{16}>1\), so
\(r=S(n^2-n)+T(n^2-1)>0\).
\end{proof}

\begin{proposition}\label{integer-reconstruction}
Every rational solution of the constructed system and the twelve
integrality interfaces has integer values in all source coordinates
apart from the positive choices that will be made for \(i,j,o,F,\tau\).
All the circuit coordinates other than \(v\) and its reciprocal-block
witnesses are integral.
\end{proposition}
\begin{proof}
The eleven ordinary inputs are integers. The guard \(\theta>0\)
forces \(b\ne0\); every power coordinate is exactly a positive power
of the integer \(b\), even across quotient steps. The remaining power
coordinates, \(\theta,e,\ell\), are therefore integers.
Lemma~\ref{lambda} gives \(\lambda\in\Z\).
All rows defining \(S,T,r,p,a,k,c,d\) and their product registers are
integer polynomial assignments. Hence they and \(\alpha,\eta\) are
integers. Lemma~\ref{rpositive} gives \(r\ge1\), so
\(c=2r+1+\varphi>0\).

Now \(vE=C_2=c^2>0\) forces \(v\ne0\). By the reciprocal
integrality interface, \(i=1/v\in\Z\setminus\{0\}\). Thus
\(E=ic^2\) and \(A_E=aE\) are integers. The norm row gives
\[
 F^2=(a^2-1)i^2c^4+1\in\Z,
\]
so \(F\in\Z\) by Lemma~\ref{monic}. Similarly,
\(\tau^2=p^2k^2-k^2+1\in\Z\) forces \(\tau\in\Z\).
The remaining definitions give integral \(G,H,I,G_H\).
\end{proof}

\begin{lemma}[The congruence quotient \(h\) is automatically positive]
\label{hpositive}
The preceding conditions force \(k>0\) and \(h\ge1\).
\end{lemma}
\begin{proof}
Here \(Y=sn^2\ge1\), \(c>0\), and
\(\eta=k^2-4(c-Yk)^2>0\). If \(k\le0\), then
\(c-Yk\ge c+|k|>|k|\), a contradiction. Thus \(k>0\).
Since \(p=2ws^2r^2n^2\ge r+2\), a negative integer \(h\)
would give \(k=r+1+h(p-1)\le0\). Therefore \(h\ge0\).
If \(h=0\), then \(k=r+1\) satisfies \(1<k<2p\). The norm row
would put the integer square \(\tau^2\) strictly between consecutive
squares:
\[
 (pk-1)^2<(pk)^2-k^2+1=\tau^2<(pk)^2.
\]
The first inequality is equivalent to \(k(2p-k)>0\); the second uses
\(k>1\). This contradiction proves \(h\ge1\).
\end{proof}

\section{Recovering positive Pell congruence witnesses}
\begin{lemma}[Pell solutions with preserved residue classes]\label{pell-lift}
Let \(D\in\Np\) be nonsquare. Suppose \(X,Y\in\Z\),
\(X>0\), \(Y\ne0\), and \(X^2-DY^2=1\). For every positive
integer modulus \(M\), there are arbitrarily large positive integers
\(X',Y'\) with
\[
 (X')^2-D(Y')^2=1,\qquad X'\equiv X\pmod M,\qquad
 Y'\equiv Y\pmod M.
\]
\end{lemma}
\begin{proof}
The unit \(U=X+Y\sqrt D\) has inverse \(X-Y\sqrt D\) in
\(\Z[\sqrt D]\). Both real conjugates are positive because
\(X>|Y|\sqrt D\). The image of \(U\) in the finite ring
\((\Z/M\Z)[T]/(T^2-D)\) is a unit, and so has a finite order \(k\).
For every integer \(t\equiv1\pmod k\), write
\(U^t=X_t+Y_t\sqrt D\). The coefficients are integers, satisfy the
same norm equation and have the required residues. If \(U>1\), take
such \(t\to+\infty\); if \(0<U<1\), take such \(t\to-\infty\).
Then \(U^t\to+\infty\), its conjugate tends to zero, and both
\(X_t,Y_t\) tend to \(+\infty\).
\end{proof}

\begin{proposition}\label{positive-reconstruction}
Every rational solution of the constructed system on the instance tail
\(x>2z\) yields a positive integer solution of all eighteen original
equations of the displayed emended source.
\end{proposition}
\begin{proof}
The integer reconstruction and previous sign arguments already give
positive values of $b,g,t,m,w,s,h,\varphi,\gamma$ and of
$\theta,\lambda,e,\ell,q,n,r,p,k,a,c,\alpha,\eta$.
In fact \(a=(wn^2+1)rsn^2>2\), and
\[
 d=(a-2)c+(4a-5)\gamma+bw>0.
\]
The source norms give \(d^2\ge a^2\) and \(F^2>1\). Replace
\(F,o\) simultaneously by \(-F,-o\) if \(F<0\); this preserves
\(I=d+oF\), every norm, and integrality. Thus assume \(F>0\).
Replace the nonzero source value \(i\) by \(|i|\), which is allowed
because the original equations contain it only through \(i^2\).

We have
\[
 G=a+F^2(d^2-a)>1,
\]
so \(D=G^2-1\) is positive and nonsquare. Let \(B_0=2r+1\).
Since \(c=B_0+\varphi>B_0>0\) and \(j\in\Z\), the integer
\(H=B_0+jc\) is nonzero. The guard \(I>0\) and the final norm
allow Lemma~\ref{pell-lift} with modulus \(Fc\).
Choose a resulting positive solution \(I',H'\) with
\(I'>d\) and \(H'>B_0\). Its residues give positive integers
\[
 o'=\frac{I'-d}{F},\qquad
 j'=\frac{H'-B_0}{c}.
\]
Replacing \(o,j\) by \(o',j'\) preserves every original equation:
they occur only in the final norm. This step proves that neither
\(o>0\) nor \(j>0\) needs its own rational positivity wrapper.
Finally replace \(\tau\) by \(|\tau|\); it is nonzero because
\(p\ge2\) and \(k>0\). All 28 original coordinates are now positive
integers and satisfy \eqref{J1}--\eqref{J18}.
\end{proof}

\section{Completeness, counts, and the final quartic}
Given a positive integer source solution, set \(v=1/i\) and assign every
power and product register its displayed value. The regrouped \(S,T\)
relations follow from \eqref{J3}. Every expression in
\eqref{guards} is a positive integer, including
\(I=d+oF>0\), so Lemma~\ref{triangular} supplies its triple.
Choose ordinary witnesses for the eleven integer interfaces, and choose
reciprocal-chart witnesses by scaling ordinary witnesses for \(i\) by
\(v\), as in Lemma~\ref{reciprocal}. This constructs a rational
solution of every new row. Together with
Proposition~\ref{positive-reconstruction}, this proves the exact
existence equivalence.

The coordinate count can be audited without expanding a high-degree
polynomial:
\begin{center}
\begin{tabular}{lr}
\toprule
Coordinates & Number\\
\midrule
The twelve free coordinates \eqref{free12}, plus \(\lambda,F,\tau\)&15\\
The fixed signed power chain &169\\
\(q_2,q_3,q_4,q_5,q_8,n,n_2\)&7\\
\(\ell,e,c_{02},c_{04},V,S,T,r,e\ell_*,g_2\)&10\\
\(s_r,s_{r2},U,p,Y,M,a,k,p_k,Y_K\)&10\\
\(d,a_c,C_2,d_2,E,A_E,F_2,G,H,I,G_H\)&11\\
\midrule
Core before positivity &222\\
Eleven positivity triples &33\\
\midrule
Rational core &255\\
Twelve integrality blocks, using existing input coordinates &\(12(m-1)\)\\
\bottomrule
\end{tabular}
\end{center}
There are 212 core rows before the positivity requirements, and 223 after
them. The twelve integrality blocks add \(12q_0\) rows. Sum the squares
of every resulting residual. All residuals have integer coefficients
and degree at most two, and at least one has a nonzero quadratic part.
The sum has degree exactly four: the sum of the squares of the
quadratic homogeneous parts cannot vanish identically over the reals.
Over \(\Q\), its vanishing is equivalent to the vanishing of every
residual. This proves Theorem~\ref{main}.

\begin{center}
\begin{tabular}{lrrrr}
\toprule
Integrality interface & \(m\)&\(q_0\)& Rational unknowns & Squares\\
\midrule
General project interface &163&137&2199&1867\\
Fully specified project interface &179&151&2391&2035\\
Ordinary-\(abc\) project interface &56&46&915&775\\
\bottomrule
\end{tabular}
\end{center}
These are baseline counts before any further algebraic sharing between
the quadratic core and the particular integrality blocks. An improvement
by such sharing must be proved for the combined system, with its
coefficient table and a complete reverse reconstruction.

\section{Verification and qualifications}
The implementation is \texttt{jones1980\_system.py}. Running
\begin{verbatim}
python jones1980_system.py --report jones1980_certificate.json
\end{verbatim}
checks the exact signed power chain, both bracket regroupings, the
source exponent in \(p\), the final squared Pell coefficient, all
quadratic degrees, all occurring coordinates, and the numerical counts.
Its structured builder exposes the eleven ordinary and one reciprocal
integer inputs separately, so a combined-system compiler can attach the
right block without silently treating \(v\) as an integer.

The program does not prove the external integrality interface. The proof
above does not identify a global minimum in the number of unknowns. The
169-step power chain is the best among the explicit window chains tested
by the accompanying search, not a proof of optimality among all
addition--subtraction chains. Nor does the cofinite argument justify
changing the exponent \(5^{60}\) to a different large exponent; the
source exponent is retained exactly.

\section{Optimization of the actual combined quadratic system}\label{sec:combined}

The project interfaces have input-square coordinates that can be reused by
the universal circuit. We now construct the combined system itself, rather
than applying the variable-count formula to independent copies and stopping
there. The transformations in this section preserve existence of rational
solutions in both directions. The arithmetic hypotheses of each interface
remain exactly those of that interface.

\subsection{An elimination valid for every quadratic integrality interface}
In the core, put \(s_r=sr\), with \(s_{r,2}=s_r^2\), and write
\(n_2=n^2\). The relevant rows are
\[
 Y=sn_2,\quad M=rY,\quad U=wn_2,\quad
 p=2U s_{r,2},\quad a=(U+1)M.
\]
The register \(M\) occurs only in its definition and in the last row.
Delete its definition and replace the last row by
\begin{equation}\label{eq:combined-a}
                  2s_r a=(p+2s_{r,2})n_2.
\end{equation}
This is an ordinary quadratic equation. The positivity argument already
proved gives \(s_r=sr>0\). Consequently the remaining rows imply
\[
 a=\frac{2(U+1)s_r^2n_2}{2s_r}=(U+1)s_rn_2,
 \qquad M=s_rn_2=rY.
\]
Thus the deleted register is recovered uniquely, the old row for \(a\)
holds, and \(a\) is an integer whenever the preceding core registers are
integers. There is no division by a possibly zero coordinate.

This alone reduces the general transfer bound from \(12m+243\) to
\(12m+242\), with \(222+12q_0\) quadratic residuals. The next three
savings use the particular integer interfaces in the project.

\begin{theorem}[Improved general quadratic transfer]
Under \eqref{integer-interface} and $\mathcal U_{\rm em}$, rational quartic
solvability is undecidable in $12m+242$ unknowns. The polynomial is a sum
of $222+12q_0$ quadratic squares.
\end{theorem}
\begin{proof}
Apply Theorem~\ref{main} and the reversible elimination
\eqref{eq:combined-a}. Every positivity and integrality argument survives
by the reconstruction just proved.
\end{proof}

\subsection{The input-square profile}
Suppose an integer block for an input \(X\) contains coordinates satisfying
\begin{equation}\label{eq:square-profile}
 X^2=uq,\qquad qq^{-1}=1,\qquad q^{-2}=(q^{-1})^2.
\end{equation}
Here \(q^{-1}\) and \(q^{-2}\) denote actual stored coordinates, with the
displayed literal equations; these symbols do not conceal extra variables.
The 163-coordinate table has
\[
 (u,q,q^{-1},q^{-2})=(Z_{37},Z_{159},Z_{160},Z_{161}).
\]
The 56-coordinate table has the simpler relation \(X^2=Z_{29}\), so all
the arguments below apply with \(u=Z_{29}\) and \(q=q^{-1}=q^{-2}=1\).
The 179-coordinate table admits the equivalent profile described below.

\subsection{Sharing the square of the integer base}
The universal power chain begins with \(b_2=b^2\). In the block testing
the integer input \(b\), write \(b^2=u_bq_b\). Delete the coordinate
\(b_2\) and its defining row. A subsequent multiplication row
\(V=b_2W\) becomes
\[
                   Vq_b^{-1}=u_bW,
\]
and a row \(V=b_2^2\) becomes
\[
                   Vq_b^{-2}=u_b^2.
\]
Both replacements are quadratic and equivalent to the original rows because
\(q_b\ne0\). Signed power-chain rows involving \(b_2\) are treated by
the same multiplication by the appropriate existing inverse coordinate.
The exact compiler checks every occurrence in the supplied chain. No other
part of the universal circuit uses \(b_2\).

\subsection{Sharing the square of \(g\)}
Originally \(g_2=g^2\) occurs only in its defining row and in the positivity
test for
\[
                   \alpha=(b-xy)q^2-e\ell g^2.
\]
The core stores a register \(e_\ell=e\ell\). Distinguish the large power
\(q=b^{5^{60}}\) in this formula from the unit \(q_g\) in the integer
block for \(g\), where \(g^2=u_gq_g\).
Replace the stored register \(e_\ell\) by \(e_\ell'=q_g e_\ell\).
Its defining equation and the expression for \(\alpha\) become
\[
 e_\ell' q_g^{-1}=e\ell,
 \qquad \alpha=(b-xy)q^2-e_\ell'u_g.
\]
These are quadratic expressions. Delete \(g_2\) and its defining row;
its value is recovered as \(u_gq_g\). The old value of \(e_\ell\) is
\(e_\ell' q_g^{-1}\). This change preserves exactly the positivity
condition on \(\alpha\), since \(e_\ell'u_g=e\ell g^2\).

\subsection{Sharing the square of the large Pell congruence modulus}
The source expression \(c=2r+1+\varphi\) is affine. Exchange the coordinate
\(\varphi\) for \(c\) by the inverse substitution
\(\varphi=c-2r-1\), and test integrality of \(c\) instead of
\(\varphi\). The value \(r\) is already integral before this step, so
the two integrality requirements are equivalent. Retain the full guard
\(c-2r-1>0\); replacing it merely by \(c>0\) would be insufficient.

In the integer block for \(c\), write \(c^2=u_cq_c\). The core register
\(C_2=c^2\) has just two further occurrences:
\[
                  d_2-a_c^2+C_2-1=0,
                  \qquad vE=C_2,
\]
where \(a_c=ac\) and \(v\) is the reciprocal input used to test
\(i=1/v\in\Z\). Replace \(C_2\) by \(u_cq_c\) in these two rows,
then delete its defining row and coordinate. Both rows remain quadratic.
The second row continues to guard \(v\ne0\): the integer block gives
\(u_cq_c=c^2\), and \(c>2r+1>0\).

The reciprocal integer block is the homogeneous block already proved:
if a source residual has homogeneous parts \(Q_2+Q_1+Q_0\), its new row is
\[
           Q_2(1,\mathbf w')+vQ_1(1,\mathbf w')+v^2Q_0=0.
\]
It equals \(v^2Q(1/v,\mathbf w'/v)\). This uses exactly \(m-1\) block
witness coordinates. The compiler checks this identity separately for
every source row, including the literal unit guards of that block.

\subsection{A three-register rechart of the 179-coordinate interface}
The existing fixed interface contains
\[
 X^2=Z_{37}Z_{163},\quad Z_{162}Z_{163}=1,
 \quad Z_{164}=Z_{163}^2.
\]
Put \(q_*=Z_{162}\), so \(Z_{163}=q_*^{-1}\) and
\(Z_{164}=q_*^{-2}\). Change precisely three coordinates:
\[
 Z_{37}^{\rm old}=q_*^2Z_{37}^{\rm new},\qquad
 Z_{77}^{\rm old}=q_*^2Z_{77}^{\rm new},\qquad
 Z_{78}^{\rm old}=q_*^4Z_{78}^{\rm new}.
\]
All other coordinates remain unchanged. The input-square row becomes
\(X^2=q_*Z_{37}^{\rm new}\). Divide the three rows named
\nolinkurl{scalar.rPair1.definition},
\nolinkurl{scalar.four_contact_product.definition}, and
\nolinkurl{scalar.forced_prime.definition} by \(q_*^4\).
Divide \nolinkurl{input.cleared.0} and \nolinkurl{input.cleared.1} by \(q_*^2\).
These are the only other rows containing a changed coordinate.

For example, the first scalar row becomes
\[
 2601(Z_{37}^{\rm new})^2
 +255102Z_{37}^{\rm new}Z_{163}
 -4Z_{77}^{\rm new}Z_{163}+6255001Z_{164}=0.
\]
Writing \(A_*\) for the fixed integer coefficient in the original
forced-normalizer row, the other two scalar rows become
\begin{align*}
 &15531948Z_{37}^{\rm new}Z_{77}^{\rm new}
 +291790992Z_{37}^{\rm new}Z_{163}
 -512(Z_{77}^{\rm new})^2\\
 &\qquad+767410608Z_{77}^{\rm new}Z_{163}
 +Z_{78}^{\rm new}Z_{90}+14315517724Z_{164}=0,\\
 &-A_*Z_{78}^{\rm new}+25Z_{78}^{\rm new}Z_{90}+Z_{164}=0.
\end{align*}
The two input-comparison rows, with \((a,b,h)=(17,23,118)\) or
\((18,24,119)\), become
\[
              (-Z_a+Z_b+2504Z_h)Z_{163}
                   +51Z_{37}^{\rm new}Z_h=0.
\]
Every displayed row is quadratic. The two literal unit rows are retained
unchanged, so the coordinate transformation is reversible on every rational
zero. This establishes the profile \eqref{eq:square-profile} with
\((u,q,q^{-1},q^{-2})=(Z_{37}^{\rm new},Z_{162},Z_{163},Z_{164})\).

\subsection{Resulting bounds and reproducible specifications}
The preceding four eliminations remove four coordinates and four rows from
the 255-coordinate, 223-row rational core. Its optimized version therefore
has 251 coordinates and 219 rows. There remain eleven ordinary integer
inputs and one reciprocal integer input. For each of the three audited
project block formats, the resulting counts are
\begin{equation}\label{eq:combined-count}
 \boxed{\quad N=251+12(m-1)=12m+239,\qquad
                R=219+12q_0.\quad}
\end{equation}

\begin{center}
\begin{tabular}{rrrr}
\toprule
Integer interface \(m\)&Block rows \(q_0\)&Rational unknowns \(N\)&Quadratic rows \(R\)\\
\midrule
163&137&2195&1863\\
179&151&2387&2031\\
56&46&911&771\\
\bottomrule
\end{tabular}
\end{center}

The parameters \(x,z,u,y\), and the fixed coefficient parameter of an
integrality interface, occur only in coefficients and are not existential
unknowns. The cofinite restriction \(x>2z\) established earlier is retained.
For the 163- and 56-coordinate families, one fixes a defining coefficient
member under the respective retained premises before applying the reduction.
The 179-coordinate interface has its explicitly specified coefficients.

The complete sparse lists of ordinary quadratic residuals are generated by
\nolinkurl{build_jones1980_combined.py}. They include all twelve actual
integer blocks, the reciprocal homogenization, and all modified core rows.
The receipts record the input source hashes, the exact variable counts,
every occurring coordinate, all rechart identities, and integrality of the
coefficient polynomials in the fixed parameters. No expanded sum of squares
is needed to specify the resulting polynomial exactly.

Finally, the defining polynomial is the sum of the squares of these
residuals. Over \(\Q\), this sum vanishes if and only if every residual
vanishes. Its degree is exactly four: at least one residual has a nonzero
homogeneous quadratic part, and a sum of real squares of such parts cannot
be the zero polynomial. The computations verify the finite algebra and the
complete specification. They do not discharge any retained arithmetic
premise, establish global minimality, or independently reprove the full
unverified universality hypothesis $\mathcal U_{\rm em}$.

\section{Source reconciliation and the remaining universality obligation}
\label{sec:source-reconciliation}

The quadratic elimination and rational-certificate arguments in this paper
are exact for the displayed polynomial system.  A separate question is
whether that precise system has the classical universality property.  The
source audit uncovered a genuine discrepancy which must be retained in
the statement of the result.

Jones's 1980 announcement, \emph{Undecidable Diophantine equations},
\emph{Bulletin of the American Mathematical Society} (New Series)
\textbf{3} (1980), 859--862, states Theorem~3 on page~861 without its
proof.  In its displayed equation for $d$, the printed coefficient is
Greek $\alpha$:
\[
                   d=bw+ca-2c+4\alpha\gamma-5\gamma.
\]
This has been checked on the original scan, by comparing the glyph with
the first-row slack $\alpha$ and with several occurrences of Latin $a$.
It is not merely an OCR discrepancy.  The construction in the present
paper instead uses
\begin{equation}\label{eq:emended-d}
                   d=bw+ca-2c+(4a-5)\gamma.
\end{equation}

An independent author-written source supports the Latin-$a$ form:
Yuri Matiyasevich, \emph{Hilbert's tenth problem: A two-way bridge between
number theory and computer science}, in C.~S.~Calude (ed.),
\emph{People and Ideas in Theoretical Computer Science}, Springer,
Singapore, 1999, pp.~177--204.  Its page~191 explicitly gives
\eqref{eq:emended-d}, in a displayed system asserted to be universal and
attributed to Jones's 1982 article.  All other equations agree with the
present system, with the following additional exception: its Pell row
is printed
\[
                         p^2k^2-k^2=1=\tau^2.
\]
The present construction uses the correct Pell equation printed in the
1980 announcement,
\begin{equation}\label{eq:emended-tau}
                         p^2k^2-k^2+1=\tau^2.
\end{equation}
The literal 1999 row cannot hold for the intended positive integers
$p\ge2$, $k\ge1$, since $(p^2-1)k^2\ge3$.  The online version also
contains a font-rendering defect in many minus signs; this is distinct
from the displayed equality-sign error.

Thus no consulted complete display is literally identical to the
system used here.  The two corrections are strongly supported by the
sources and by the local mathematics, but that evidence is not a
replacement for checking the full universality proof.  The following
elementary computation explains the role of Latin $a$.

\begin{lemma}[The power-of-two Pell congruence]
Let $a\ge2$ be an integer and define integers $X_n,Y_n$ by
\[
 X_n+Y_n\sqrt{a^2-1}=(a+\sqrt{a^2-1})^n
 \qquad(n\ge0).
\]
Then
\[
             X_n-(a-2)Y_n\equiv2^n\pmod{4a-5}.
\]
\end{lemma}
\begin{proof}
Set $R_n=X_n-(a-2)Y_n$.  The two standard multiplication recurrences
give
\[
 R_{n+1}=2X_n+(2a-1)Y_n,
 \qquad
 R_{n+1}-2R_n=(4a-5)Y_n.
\]
Since $R_0=1$, induction proves the congruence.  No relation with the
unrelated slack $\alpha$ occurs in this calculation.
\end{proof}

The 1975 paper of Matiyasevich and Robinson,
\emph{Reduction of an arbitrary Diophantine equation to one in
13 unknowns}, \emph{Acta Arithmetica} \textbf{27} (1975), 521--553,
Theorem~4 on pages~532--533, independently supplies the classical
three-Pell-equation bridge.  Its notation and its auxiliary factors
differ from the system above.  Invoking that theorem alone does not
prove every specialization and substitution needed for the complete
Jones system, and in particular does not settle the above source
discrepancy by itself.

For a fully rigorous checkpoint the extra premise is therefore stated
explicitly:
\begin{quote}
\textbf{Hypothesis $\mathcal U_{\rm em}$:}
The displayed 28-variable positive-integer system, with
\eqref{eq:emended-d} and \eqref{eq:emended-tau}, has the universality
property asserted in Jones's theorem; in particular, for some fixed
positive integer coefficients $z,u,y$, its solvable positive input
values $x$ form a nonrecursive set.
\end{quote}
Under this premise, the exact quadratic reduction and its stated
variable count give the claimed undecidability transfer.  Without it,
the coefficient construction and both directions of its equivalence
with the displayed system remain proved, but the undecidability label
is not yet established.  This is a source-verification premise and is
unrelated to $abc$.

The remaining task is concrete: obtain and verify the relevant complete
proof in Jones, \emph{Universal diophantine equation},
\emph{Journal of Symbolic Logic} \textbf{47} (1982), 549--571,
DOI~10.2307/2273588, or prove directly that the precise emended system is
equivalent to an independently verified universal exponential or
binomial-coefficient system.  It is not enough to check the local Pell
congruence, or to copy an unverified secondary transcription.

The original scan used for the 1980 comparison is reproduced on PDF
page~53 (zero-based page~52) of Gabriel Robins's University of Virginia
slides:
\url{https://www.cs.virginia.edu/~robins/cs3102/slides/Theory_Slides_Russell_Hardy_Ramanujan_Ramsey_Hilbert_v61.pdf}.
The author-written 1999 chapter is available from the editor's academic
site at
\url{https://www.cs.auckland.ac.nz/~cristian/tcspi/ch9.pdf}.
The 1975 authors' scan is available at
\url{https://logic.pdmi.ras.ru/~yumat/papers/23_paper/23_home.html}.

\begin{thebibliography}{9}
\bibitem{Mat99} Yu. Matiyasevich,
\emph{Hilbert's tenth problem: A two-way bridge between number theory and computer science},
in C.~S.~Calude (ed.), \emph{People and Ideas in Theoretical Computer Science},
Springer, Singapore, 1999, pp.~177--204.
\url{https://www.cs.auckland.ac.nz/~cristian/tcspi/ch9.pdf}.
\bibitem{Jones80} J. P. Jones,
\emph{Undecidable Diophantine equations},
Bulletin of the American Mathematical Society (New Series)
\textbf{3} (1980), 859--862.
\href{https://doi.org/10.1090/S0273-0979-1980-14832-6}{doi:10.1090/S0273-0979-1980-14832-6}.
Theorem~3, p.~861; its printed display is compared with the emended system above.
\bibitem{Jones82} J. P. Jones,
\emph{Universal diophantine equation}, Journal of Symbolic Logic
\textbf{47} (1982), 549--571.
\href{https://doi.org/10.2307/2273588}{doi:10.2307/2273588}.
\bibitem{MT} A. Mishchenko and A. Treier,
\emph{Universal input of knapsack problem for groups},
Journal of Physics: Conference Series \textbf{1050} (2018), 012054.
\href{https://doi.org/10.1088/1742-6596/1050/1/012054}{doi:10.1088/1742-6596/1050/1/012054}.
Its explicit transcription is not used for the present theorem.
\bibitem{MR} Yu. Matiyasevich and J. Robinson,
\emph{Reduction of an arbitrary Diophantine equation to one in 13
unknowns}, Acta Arithmetica \textbf{27} (1975), 521--553.
The author's public copy is at
\url{https://logic.pdmi.ras.ru/~yumat/papers/23_paper/23_home.html}.
\bibitem{Scan} G. Robins, University of Virginia, theory lecture slides,
\emph{Russell, Hardy, Ramanujan, Ramsey, Hilbert}, version 61,
PDF page 53 reproduces the original Jones page 861.
\url{https://www.cs.virginia.edu/~robins/cs3102/slides/Theory_Slides_Russell_Hardy_Ramanujan_Ramsey_Hilbert_v61.pdf}.
Used to inspect the primary-source mathematical display, not as a
replacement theorem.
\end{thebibliography}
\end{document}
